\scanpage{7}
\mspage{3}
% archive sheet number: 6
\deleted{\unsure{221}}

$\deleted{X\mathrel{>\mkern-12mu<}Y}\Longleftrightarrow$

\deleted{\unsure{donc}} Appliquons la définition générale de $\mathrel{>\mkern-12mu<}$,

\deleted{\unsure{On voit} que pour $X\ne Y$, on a $X$}

\deleted{On voit que pour $X\ne Y$, on ne peut avoir $\widetilde X\cap\widetilde Y\ne\varnothing$ que si $\unsure{X}\ \delta$}
\[
X\ne Y\Longrightarrow\widetilde X\cap\widetilde Y=\delta X\cap\delta Y\deleted{\unsure{\text{est}}},\ \text{de cardinal }0\text{ ou }1
\]
\deleted{et si $\widetilde X\cap\widetilde Y\ne\varnothing$, on est dans la situation}

$\deleted{X\triangleleft Y\text{ ou }Y\triangleleft X\text{ ou}}$

$\deleted{X,Y\in\unsure{\mathscr M|\mathcal M}_1}$\quad $\deleted{\unsure{\delta}}\operatorname{card}(\delta X\cap\delta Y)=1$

\deleted{Donc $X\mathrel{>\mkern-12mu<}Y\Longleftrightarrow X=Y$ \unsure{et} ou $X\triangleleft Y$ ou $Y\triangleleft X$}

\uline{Proposition}\quad $X\mathrel{>\mkern-12mu<}Y\Longleftrightarrow X=Y\text{ ou }X|\circ|Y$ i.e. $\{X,Y\}\in\operatorname{Drap}(\mathscr M)$
\deleted{\unsure{ou}}

\deleted{\unsure{Dém}} \uline{Dém.}\quad \unsure{ONS|OPS} $X\ne Y$. Alors $\deleted{\unsure{\widetilde X}}$
\[
\widetilde X\cap\widetilde Y=\delta X\cap\delta Y\quad\text{est de cardinal }0\text{ ou }1
\]
distinguons les deux cas.

$\deleted{\delta X\cap\delta Y=\widetilde X\cap\widetilde Y}$

a) $\widetilde X\cap\widetilde Y=\varnothing$. Alors la définition \inserted{de $X\mathrel{>\mkern-12mu<}Y$} signifie $\widetilde X|\bullet|\widetilde Y$, et on voit que cela équivaut aussi à $X|\circ|Y$, par l'hypothèse $\deleted{\unsure{\delta}}\widetilde X\cap\widetilde Y=\varnothing$.

b) $\operatorname{card}\widetilde X\cap\widetilde Y=\deleted{\varnothing}1$, i.e. $\widetilde X\cap\widetilde Y=\{b\}$. \deleted{Alors} \deleted{\unsure{Deux}}

\deleted{$X\mathrel{>\mkern-12mu<}Y$ signifie} Si $X\triangleleft Y$ ou $Y\triangleleft X$, alors $X\mathrel{>\mkern-12mu<}Y$ et $X|\circ|Y$, la prop. est satisfaite. OR \deleted{\unsure{que}} \deleted{\unsure{X} que $X$ et $Y$ \unsure{ne sont}} que l'on n'est pas dans ce cas, mais alors $X,Y\in\mathscr M_1$, \inserted{$X,Y\notin\widetilde X\cap\widetilde Y$,} \deleted{et \unsure{nous avons}} $X\mathrel{>\mkern-12mu<}Y$ \deleted{signifie} implique $X|\circ|Y$, \deleted{\unsure{justement}}, donc (quitte à intervertir)
\[
X=I_{a,b},\quad Y=I_{b,c},\quad a<b<c.
\]
Mais alors $(\widetilde X\setminus\{b\})|\circ|(\widetilde Y\setminus\{b\})$ i.e. $X\mathrel{>\mkern-12mu<}Y$, qed.

\uline{Corollaire} Pour qu'une partie \inserted{finie} $\unsure{\Phi|\mathscr F}$ de $\mathscr M$ soit une préfigure de $\mathscr M$ (i.e. $\overline\Phi$ soit une figure) il f. et s. qu'elle soit totalement ordonnée par $\leq$.
